\section{公理集合の変更に伴う $W_K$ の変化}
\label{sec:WK-set-change}

本節では、前節までに定義した $W_K$ について、比較対象となる公理集合に公理を追加・削除した場合の変化を整理する。以下では特記しない限り、Solver集合 $\mathcal{Z}$、配置値 $x_z(a)$、したがって既存公理間の $P$、$D_{\mathrm{old}}$、$Q$ を固定する。これは比較対象集合だけを変更する制御検証である。公理の追加がSolver集合や既存公理の配置自体を変える場合には、その変化を別途扱わなければならない。

\subsection{一つの公理を追加する場合}

変更前の公理集合を $\mathcal{A}_n$、$|\mathcal{A}_n|=n>1$ とし、対象公理 $a\in\mathcal{A}_n$ とする。新しい公理 $c\notin\mathcal{A}_n$ を追加し、$\mathcal{A}_{n+1}=\mathcal{A}_n\cup\{c\}$ とする。既存の比較差分の合計を
\[
S_a=\sum_{\substack{b\in\mathcal{A}_n\\b\ne a}}D_{\mathrm{old}}(a\mid b)
\]
と置くと、追加前後は
\[
W_n(a)=P(a)+\frac{S_a}{n-1},\qquad
W_{n+1}(a)=P(a)+\frac{S_a+D_{\mathrm{old}}(a\mid c)}{n}
\]
である。従って
\[
\boxed{
W_{n+1}(a)-W_n(a)
=\frac{D_{\mathrm{old}}(a\mid c)-M_n(a)}{n}
}
\tag{5}
\]
を得る。ここで $M_n(a)=S_a/(n-1)$ は追加前の比較差分の平均である。

式 (5) より、新公理との方向付き差分が既存平均より大きければ $W_K(a)$ は増加し、小さければ減少し、等しければ不変である。したがって、$P(a)$ が固定されていても $W_K(a)$ は不変とは限らない。$n=1$ の場合は追加前に比較対象がないため、取り決め $W_1(a)=P(a)$ を用いると
\[
W_2(a)-W_1(a)=D_{\mathrm{old}}(a\mid c)
\]
であり、式 (5) の平均表現は適用しない。

\subsection{$Q$ による同値な表現}

$D_{\mathrm{old}}(a\mid b)=P(a)-Q(a,b)$ を用い、既存の積平均を
\[
\overline Q_n(a)=\frac{1}{n-1}\sum_{\substack{b\in\mathcal{A}_n\\b\ne a}}Q(a,b)
\]
とすれば、式 (5) は
\[
\boxed{
W_{n+1}(a)-W_n(a)
=\frac{\overline Q_n(a)-Q(a,c)}{n}
}
\tag{6}
\]
となる。従って、$W_K(a)$ の増減は、新公理との積平均が既存の積平均より小さいか、大きいかによっても判定できる。これは式 (5) と同一の結論を異なる形で表したものである。

\subsection{複数の公理をまとめて追加する場合}

追加する公理集合を $\mathcal{C}$ とし、$|\mathcal{C}|=m\ge1$、$\mathcal{C}\cap\mathcal{A}_n=\varnothing$ とする。対象公理 $a\in\mathcal{A}_n$ について、追加公理との平均差分を
\[
M_{\mathcal{C}}(a)=\frac1m\sum_{c\in\mathcal{C}}D_{\mathrm{old}}(a\mid c)
\]
と定義する。$n>1$ ならば
\[
\boxed{
W_{n+m}(a)-W_n(a)
=\frac{m}{n+m-1}
\bigl(M_{\mathcal{C}}(a)-M_n(a)\bigr)
}
\tag{7}
\]
である。変化方向は、追加公理群との平均差分と既存平均の比較によって決まる。追加数 $m$ は変化量の係数に影響するが、平均差の符号が変化方向を決める。

\subsection{一つの公理を削除する場合}

変更前の集合を $\mathcal{A}_n$、$n>2$ とし、削除対象を $c$、残す対象公理を $a\ne c$ とする。削除前の平均を $M_n(a)$ とすれば、削除後の比較平均は削除前の全体平均から一項を除いた平均となる。よって
\[
\boxed{
W_{n-1}(a)-W_n(a)
=\frac{M_n(a)-D_{\mathrm{old}}(a\mid c)}{n-2}
}
\tag{8}
\]
である。削除された公理との差分が削除前の平均より小さい場合、残った比較差分の平均は上がり、$W_K(a)$ は増加する。削除対象との差分が平均より大きい場合は減少し、等しい場合は不変である。

削除後に公理が一つだけ残る場合は $W_1(a)=P(a)$ とする。比較対象が空になるため、式 (8) の分母 $n-2$ を用いる一般式は適用しない。

\subsection{複数の公理をまとめて削除する場合}

変更前の公理集合を互いに素な集合 $\mathcal{R}$ と $\mathcal{C}$ の和集合
\[
\mathcal{A}_n=\mathcal{R}\cup\mathcal{C},\qquad
\mathcal{R}\cap\mathcal{C}=\varnothing
\]
とし、$|\mathcal{R}|=r>1$、$|\mathcal{C}|=m\ge1$ とする。対象公理 $a$ は $\mathcal{R}$ に属する。残す公理群と削除公理群への平均差分をそれぞれ
\[
M_{\mathcal{R}}(a)=\frac1{r-1}\sum_{\substack{b\in\mathcal{R}\\b\ne a}}D_{\mathrm{old}}(a\mid b),\qquad
M_{\mathcal{C}}(a)=\frac1m\sum_{c\in\mathcal{C}}D_{\mathrm{old}}(a\mid c)
\]
とする。すると
\[
\boxed{
W_{\mathcal{R}}(a)-W_{\mathcal{A}_n}(a)
=\frac{m}{r+m-1}
\bigl(M_{\mathcal{R}}(a)-M_{\mathcal{C}}(a)\bigr)
}
\tag{9}
\]
となる。これは複数公理の追加式 (7) と同じ加重平均の構造を持つ。

\subsection{追加と削除の対称性、追加順序}

Solver集合と配置値を固定したまま、同じ最終公理集合を得るなら、$W_K(a)$ は最終集合と各方向付き差分だけで決まるため、追加順序に依存しない。逐次追加時の各段階の変化量は、その時点の比較対象数と平均が変わるため異なり得るが、
\[
\sum_{j=0}^{m-1}\bigl(W_{n+j+1}(a)-W_{n+j}(a)\bigr)
=W_{n+m}(a)-W_n(a)
\]
であり、最終的な差は同じである。

同様に、追加と削除を逆向きに行って元の公理集合に戻れば、固定条件のもとで最終的な $W_K$ も元に戻る。ただし、追加時と削除時の各段階の変化量が同じ数値になるわけではない。

\section{正規化後の $K$ の変化}

正規化前の最大値を
\[
H_n=\max_{b\in\mathcal{A}_n}W_n(b)
\]
とし、$H_n>0$ のとき $K_n(a)=W_n(a)/H_n$ とする。公理追加後も $H_{n+1}>0$ であると仮定する。このとき
\[
K_{n+1}(a)=K_n(a)
\quad\Longleftrightarrow\quad
W_{n+1}(a)H_n=W_n(a)H_{n+1}.
\tag{10}
\]
特に $W_n(a)>0$ ならば、$K$ が不変である必要十分条件は
\[
\frac{W_{n+1}(a)}{W_n(a)}=\frac{H_{n+1}}{H_n}
\]
である。つまり、分子と正規化分母が同じ比率で変化するときに限り、正規化後の値が不変となる。

$\Delta W(a)=W_{n+1}(a)-W_n(a)$、$\Delta H=H_{n+1}-H_n$ とすれば
\[
\boxed{
K_{n+1}(a)-K_n(a)
=\frac{H_n\Delta W(a)-W_n(a)\Delta H}
{H_n(H_n+\Delta H)}
}
\tag{11}
\]
である。分母は正なので、増減は分子の符号で判定される。従って $W_K(a)$ の変化だけでは正規化後の $K(a)$ の変化は決まらず、全公理の最大値の変化も必要である。

なお、最大値 $H_n$ は、公理追加時に必ず増加するとは限らない。既存公理の $W_K$ が減少し得るためである。また、最大値が $0$ であれば正規化は未定義であり、別の規則を設けない限り式 (10)--(11) は適用できない。

\section{既存公理間の順位と逆転条件}

同一の公理集合では、正の共通分母による正規化は順位を変えない。従って、正規化後の $K$ の順位は $W_K$ の順位と一致する。

既存公理 $a,b$ について、公理追加前後の順位差は
\[
W_{n+1}(a)-W_{n+1}(b)
=\bigl(W_n(a)-W_n(b)\bigr)
+\bigl(\Delta W(a)-\Delta W(b)\bigr)
\tag{12}
\]
である。追加前に $W_n(a)>W_n(b)$ であった場合、順位が逆転する必要十分条件は
\[
\boxed{\Delta W(b)-\Delta W(a)>W_n(a)-W_n(b)}
\tag{13}
\]
である。等号なら追加後に同順位となる。

各差分は $[0,1]$ に属するため、$n>1$ の一公理追加について
\[
-\frac1n\le\Delta W(a)\le\frac1n
\]
であり、
\[
-\frac2n\le\Delta W(a)-\Delta W(b)\le\frac2n
\]
を得る。従って、追加前の正の順位差が $2/n$ 以上なら、一回の追加で順位逆転することはできない。この条件は逆転不能を示す十分な上限条件であり、順位差が $2/n$ 未満であること自体は逆転を保証しない。実際の判定には式 (13) が必要である。

\section{確定事項・適用条件・保留事項}
\begin{enumerate}
\item 追加・削除による $W_K$ の変化は、変更される比較対象群と維持される比較対象群の平均差分の差として表せる。
\item 追加後の個別の $W_K$ は変化し得るが、Solver集合・配置値を固定して同じ最終公理集合に至るなら、最終値は追加順序に依存しない。
\item 正規化後の $K$ は、個別の $W_K$ の変化だけでなく、全公理の最大値の変化にも依存する。
\item 既存公理間の順位逆転は、個別の変化量の差が追加前の順位差を超えるときに生じる。
\item 以上の結論は、有限な公理集合、有限で空でないSolver集合、等重み平均、および既存の配置・関係値を固定する制御条件のもとで成立する。
\item 公理の変更がSolver集合や既存公理の配置値そのものを変える場合には、その効果を今回の比較対象数の変化による効果と分離して評価する必要がある。
\item 連続値 $x_z(a)\in[0,1]$ のECA上の意味、配置値の正当化、非有限Solver集合・非可測集合への拡張は、本節の代数的結果だけでは確定しない。
\end{enumerate}
